\clearpage\thispagestyle{empty}
\pdfbookmark[0]{Sources for the March 1805 Ceres Report}{ceres180503-sources}
\begin{center}\textbf{Sources for the March 1805 Ceres Report}\end{center}
\noindent This AI translation follows Franz Xaver von Zach’s report as excerpted in \href{https://gdz.sub.uni-goettingen.de/id/PPN236007467}{Gauss’s \textit{Werke}, vol.~VI} (Göttingen, 1874), pp.~261--262. The corresponding passages of \textit{Ueber die Ceres}, article XXVI of the \textit{Monatliche Correspondenz}, vol.~XI, March 1805, are on \href{https://www.digitale-sammlungen.de/en/view/bsb10538603?page=295}{p.~283} and \href{https://www.digitale-sammlungen.de/en/view/bsb10538603?page=298}{pp.~286}--\href{https://www.digitale-sammlungen.de/en/view/bsb10538603?page=299}{287}.
\par\smallskip\noindent Carl Friedrich Gauss calculated the tenth set of elements. Friedrich Wilhelm Bessel supplied the positions of the Sun needed for the ephemeris. Von Zach refers to Bessel’s \textit{Berechnung der Harriot’schen und Torporley’schen Beobachtungen des Cometen von 1607}, in the November 1804 issue, \href{https://www.digitale-sammlungen.de/en/view/bsb10538602?page=439}{p.~425} and \href{https://www.digitale-sammlungen.de/en/view/bsb10538602?page=440}{following}.
\par\smallskip\noindent Translation and typesetting produced by AI. Journal scans from the Bayerische Staatsbibliothek.
\clearpage
